% Generated by code/tables/tabA-measurement.py. Do not edit by hand.

\begin{table}[htbp]
\centering
\begin{threeparttable}
\caption{The subtask decomposition: counts, time shares, and activity types}
\label{tab:subtasks}
\footnotesize
\begin{spacing}{1}
\begin{tabular}{@{}ll@{}}
\toprule
 &  \\
\midrule
\multicolumn{2}{@{}l}{\emph{A. Subtasks per task}} \\
3 subtasks & 3 tasks (0.0 percent) \\
4 subtasks & 2,001 tasks (10.4 percent) \\
5 subtasks & 12,518 tasks (64.9 percent) \\
6 subtasks & 4,759 tasks (24.7 percent) \\
Total & 99,157 subtasks in 19,281 tasks, 5.14 per task \\
Time share of a subtask & median 15 percent, 10th to 90th percentile 10 to 35 \\
\addlinespace
\multicolumn{2}{@{}l}{\emph{B. Share of task time by activity type (percent)}} \\
Physical execution & 24.6 \\
Gathering information & 20.1 \\
Processing information & 18.3 \\
Communication and coordination & 12.9 \\
Administrative and procedural work & 10.0 \\
Generating content & 8.0 \\
Decision and judgment & 6.1 \\
\bottomrule
\end{tabular}
\end{spacing}
\begin{tablenotes}\notesize
\item Notes: Prompt 0 (Appendix~\ref{app:prompts}) split each of the 19,281 O*NET task statements into three to seven subtasks with integer time shares summing to 100 and one of seven activity types; no task received seven. Panel B weights each subtask by its time share.
\end{tablenotes}
\end{threeparttable}
\end{table}
